\section{Auxiliary primes and primary generators}
\label{sec:table}

The splitting procedure needs explicit generators of small-prime parts of
multiplicative groups.  This section constructs those generators from the
auxiliary primes supplied in Theorem~\ref{thm:reduction}.
The construction uses linear algebra
in cyclotomic algebras and calls the factorization procedure only in strictly
smaller degrees.  Existence of small auxiliary primes is a separate
number-theoretic question; no size bound for them is assumed in this section.

We use the notation $v_q(a)$ and the full $q$-primary subgroup
introduced in Section~\ref{ff:section}.

Throughout the construction, $p>n\ge2$.  For each prime $q\le n$,
recall that the supplied auxiliary prime $\ell_q$ satisfies
\begin{equation}\label{table:primes}
 \ell_q\notin\{2,3,p,q\},\qquad
 \ell_q\equiv1\pmod{12q},\qquad
 p^{(\ell_q-1)/q}\not\equiv1\pmod{\ell_q}.
\end{equation}
Recall the bound $E=\max(2,\max_{q\le n}\ell_q)$ from
Theorem~\ref{thm:reduction}.  Supplied primes can be checked by trial
division and modular exponentiation in time polynomial in $E$, $n$, and
$\log p$.  The same tests can search upward for the first suitable prime;
if one exists below $E$, that search also takes time polynomial in these
parameters.  Factoring $\ell_q-1$ is unnecessary.

The table has the following contents:
\begin{itemize}
\item for $q=2$, an element $\omega_2\in\mathbb F_p^*$ of order
      $2^{v_2(p-1)}$;
\item for each odd prime $q\le n$, an explicitly represented field
      $K_q/\mathbb F_p$, an element $\zeta_q\in K_q$ of order $q$ with
      $K_q=\mathbb F_p(\zeta_q)$, and an element $\omega_q\in K_q^*$
      of order $q^{v_q(|K_q|-1)}$.
\end{itemize}
An explicit field here consists of an $\mathbb F_p$-basis, its
multiplication constants, and the coordinates of its identity.  A quotient
by a specified irreducible polynomial is an equivalent representation.
Both descriptions support arithmetic and linear algebra in polynomial time.
We use the standard structure theory of finite fields; see
\cite{LidlNiederreiter97}.

\subsection{A degree-\texorpdfstring{$q$}{q} field without factoring a cyclotomic polynomial}

Fix $q$ and $\ell=\ell_q$.  Put
\[
 R_\ell=\mathbb F_p[T]/(1+T+\cdots+T^{\ell-1}),\qquad
 \Gamma_\ell=(\mathbb Z/\ell\mathbb Z)^*,\qquad
 H_q=\{a^q:a\in\Gamma_\ell\}.
\]
For $a\in\Gamma_\ell$, let $\tau_a$ be the automorphism of
$R_\ell$ defined by substitution $T\mapsto T^a$.  Enumerate $H_q$ by taking $q$th powers modulo $\ell$,
and compute
\begin{equation}\label{table:invariants}
 U_q=R_\ell^{H_q}
    =\bigcap_{a\in H_q}\ker(\tau_a-\mathrm{id})
\end{equation}
by $\mathbb F_p$-linear algebra.  Products of basis
elements of this space, reduced in $R_\ell$ and expressed in its computed
basis, give the multiplication constants of $U_q$.

\begin{lemma}\label{table:degreeq}
The algebra $U_q$ is a field of degree $q$ over $\mathbb F_p$.
\end{lemma}
\begin{proof}
Since $p\ne\ell$, the polynomial defining $R_\ell$ has $\ell-1$
distinct roots over an algebraic closure $\overline{\mathbb F}_p$.
They are precisely the primitive $\ell$th roots of unity.  Evaluation at
these roots identifies the scalar extension of $R_\ell$ with the algebra
of functions on that set.  The group $\Gamma_\ell$ acts regularly on the
set, and its subgroup $H_q$ has index $q$.  Taking kernels commutes with
extension of scalars.  Consequently
\[
 U_q\otimes_{\mathbb F_p}\overline{\mathbb F}_p
       \simeq\overline{\mathbb F}_p^{\,q},
\]
with the factors indexed by the $H_q$-orbits.

Frobenius permutes the primitive roots by raising them to the $p$th
power, hence acts on the orbit set by multiplication by the class of $p$
in $\Gamma_\ell/H_q$.  A cyclic group of order $\ell-1$ has
$H_q=\{a:a^{(\ell-1)/q}=1\}$, so the last condition in
\eqref{table:primes} says that this class is nonidentity.  The quotient has
prime order $q$; its nonidentity elements act transitively on it.
A finite reduced algebra over a finite field is a product of finite fields,
and its field factors correspond exactly to Frobenius orbits on its
geometric points.  There is one such orbit, of length $q$, proving the
claim.
\end{proof}

This construction is useful precisely because the computation of
$U_q$ never asks for the individual irreducible factors of the degree
$\ell-1$ cyclotomic polynomial.  Its field property follows from the
verified power-residue test on the integer prime $\ell$.

\subsection{Extracting the primary generators}

For $q=2$, solve the $\mathbb F_p$-linear equation
\[
 v^p=-v\qquad(v\in U_2)
\]
and take a nonzero solution.  The trace map $v\mapsto v+v^p$ from the
quadratic field $U_2$ to $\mathbb F_p$ is nonzero, because it sends $1$
to $2$.  Its kernel has dimension one, so this instruction succeeds.
Then $c=v^2$ lies in $\mathbb F_p^*$ and
\[
 c^{(p-1)/2}=v^{p-1}=-1.
\]
Thus $c$ is a nonsquare.  With $s=v_2(p-1)$, set
\begin{equation}\label{table:quadratic}
 \omega_2=c^{(p-1)/2^s}.
\end{equation}
The order of a nonsquare in a cyclic group of even order has the full
$2$-part of that order.  The odd exponent in \eqref{table:quadratic}
therefore leaves $\omega_2$ of exact order $2^s$.

For odd $q$, first factor
\[
 \Phi_q(Z)=1+Z+\cdots+Z^{q-1}
\]
over $\mathbb F_p$.  The smaller-degree nature of this call will be
justified below.  Choose the first monic irreducible factor in a fixed
coefficient order, form its quotient field $K_q$, and let $\zeta_q$ be the
image of $Z$.  Every root of $\Phi_q$ has order $q$, because $p\ne q$.
Writing $d_q=[K_q:\mathbb F_p]$, finite-field theory gives
\[
 d_q=\operatorname{ord}_q(p)\mid q-1,
\]
where $\operatorname{ord}_q(p)$ is the multiplicative order of $p$
modulo $q$.  In particular $\gcd(d_q,q)=1$.  The tensor product
\[
 V_q=U_q\otimes_{\mathbb F_p}K_q
\]
is therefore a field of degree $q$ over $K_q$: finite fields of coprime
degrees intersect in $\mathbb F_p$, so they are linearly disjoint.
Its multiplication table is obtained directly from those of its two
factors.

Set $Q_q=|K_q|$ and $s_q=v_q(Q_q-1)$.  Solve the $K_q$-linear system
\begin{equation}\label{table:eigenvector}
 u^{Q_q}=\zeta_q u\qquad(u\in V_q)
\end{equation}
and choose a nonzero solution.  To prove existence, the homomorphism
$x\mapsto x^{Q_q-1}$ on $V_q^*$ has image equal to the norm-one subgroup
for $V_q/K_q$.  Indeed both subgroups have order
$(Q_q^q-1)/(Q_q-1)$.  The norm of $\zeta_q$ equals
$\zeta_q^{1+Q_q+\cdots+Q_q^{q-1}}=\zeta_q^q=1$.
Thus $\zeta_q$ lies in this image, proving that
\eqref{table:eigenvector} has a nonzero solution.

For odd $q$ and $Q_q\equiv1\pmod q$, the lifting-the-exponent identity is
\begin{equation}\label{table:lte}
 v_q(Q_q^a-1)=s_q+v_q(a)\qquad(a\ge1).
\end{equation}
This is Equation~\eqref{ff:lte}.  In particular $v_q(Q_q^q-1)=s_q+1$.  Equation
\eqref{table:eigenvector} says that $u^{Q_q-1}$ has order exactly $q$,
so the $q$-part of the order of $u$ is exactly $q^{s_q+1}$.  Consequently
\[
 w=u^{(Q_q^q-1)/q^{s_q+1}},\qquad \omega_q=w^q
\]
have orders $q^{s_q+1}$ and $q^{s_q}$, respectively.  The unique
subgroup of $V_q^*$ of order $q^{s_q}$ lies in $K_q^*$, since
$q^{s_q}\mid Q_q-1$.  Therefore $\omega_q\in K_q$, as required.
Its coordinates in $K_q$ are obtained by solving the linear equations for
the known inclusion $K_q\hookrightarrow V_q$.

All integer valuations used here are found by repeated division by the
specified prime $q$.  All powers use binary exponentiation.  Neither
construction requires an integer factorization of $p-1$ or $Q_q-1$.

\subsection{The increasing-degree construction}

The only call above to an as yet unspecified factoring procedure is the
factorization of $\Phi_q$, whose degree is $q-1$.  The interface with the
rest of the paper is the following.

\begin{proposition}[Table construction]\label{table:construction}
Suppose the auxiliary primes \eqref{table:primes} are supplied, and
suppose a deterministic procedure has the following uniform contract:
for every integer $1\le b<n$, it factors every monic polynomial of degree at
most $b$ over $\mathbb F_p$ using table entries only for primes at most
$b$.  Then the complete table through $n$ can be built in increasing
prime order.  Its construction calls that procedure at most once for each
odd prime $q\le n$, on a polynomial of degree $q-1$; each such call uses
only already constructed table entries.

Apart from these factorization calls, the construction has bit cost
polynomial in $E$, $n$, and $\log p$, with an absolute exponent.  Each
stored odd-prime field has degree at most $q-1$, and the temporary field
$V_q$ has degree at most $q(q-1)$ over $\mathbb F_p$.
\end{proposition}
\begin{proof}
Begin with the quadratic construction, which needs no factorization call.
For an odd prime $q$, every prime at most $q-1$ is smaller than $q$, so
the indicated factorization call has all its required entries.  Its output
constructs $K_q$, and the preceding linear-algebra construction supplies
$\omega_q$.  This proves the inductive assertion and the degree bounds.

For completeness, the largest auxiliary algebra $R_{\ell_q}$ has
dimension $\ell_q-1\le E$.  Substitution operators are computed by
modular polynomial powering.  There are at most $\ell_q-1$ elements in
$H_q$, and their fixed-space equations have at most $\ell_q-1$
unknowns.  Gaussian elimination gives its invariant subspace and its
multiplication constants in polynomial time.  All subsequent field
dimensions are at most $n^2$.  Including the modular residue tests, the exponent integers used have
$O(\log(E+1)+n^2\log p)$ bits, and each extraction of a specified prime valuation
has polynomial cost.  Schoolbook arithmetic consequently gives, for
example, an upper bound
\[
 O\bigl(n(E+1)^{10}(n+\lceil\log_2p\rceil+1)^{40}\bigr)
\]
for the work outside the factorization calls.  The same generous bound
includes trial division and upward searches through $E$.
\end{proof}

Section~\ref{driver:section} proves that \textsc{Factor} has exactly this
interface.  Its correctness and the table construction are established
by a common induction: the entry for $q$ invokes factorization only below
$q$, whereas factorization in degree $b$ invokes entries only at most
$b$.  Thus the table uses no factorization oracle.  The only external
data in this construction are the verified auxiliary primes.
